\documentclass[11pt,a4paper]{article}

\usepackage{Act}
\usepackage{qrcode}

\begin{document}
\input{\detokenize{/home/fenarius/Travail/Cours/cpge-info/latex/Macros.tex}}

% Entête de la fiche (à modifier avec la macro correspondante dans le fichier macros)
\Fiche{Énigmes de programmation}{\textsc{itc}}

\pythonmode

On dit qu'un entier $n > 0$ est \textit{abondant} lorsqu'il est strictement inférieur à la somme de ses diviseurs stricts positifs (c'est-à-dire strictement inférieurs à $n$). Par exemple :
\begin{itemize}
 \item [\textbullet] Les diviseurs stricts de $20$ sont $1, 2, 4, 5$ et $10$, de somme $23$, donc $20$ est abondant.
 \item [\textbullet] Les diviseurs stricts de $32$ sont $1, 2, 4, 8$ et $16$, de somme $31$, donc $32$ n'est pas abondant.
 \item [\textbullet] Les diviseurs stricts de $28$ sont $1, 2, 4, 7$ et $14$, de somme $28$, donc $28$ n'est pas abondant ($28$ est dit \textit{parfait}).
\end{itemize}

\begin{tcolorbox}[left=0cm,title=\textbf{\faQuestion{} Nombres abondants},colbacktitle=cfond]
    Combien y a-t-il de nombres abondants strictement inférieurs à $100\,000$ ?
\end{tcolorbox}

\end{document}
